% GENERATED FILE. Edit canonical lessons, then run npm run generate:books.
% canonical-source-hash: fnv1a64:06ffbd0d307af642
% canonical-lessons: FA-C266-chaharrah, FA-C266-mahalle, FA-C266-sakhteman, FA-C266-vorudi, FA-C266-khoruji

\chapter{Crossroads, Neighbourhood, Building, Entrance, Exit}
\label{ch:fa-crossroads-neighbourhood-building-entrance-exit}
\hlchaptermodality{266}

\begin{chapteropening}
\textbf{By the end of this chapter:} \emph{I can say a crossroads, a neighbourhood, a building, an entrance and an exit.}

Complete the last lesson of chapter 266: I can say a crossroads, a neighbourhood, a building, an entrance and an exit.
\end{chapteropening}

\section[\texorpdfstring{\fa{چهارراه} (chahârrâh)}{چهارراه (chahârrâh)}]{\fa{چهارراه} (chahârrâh) — a crossroads}

\label{lesson:FA-C266-chaharrah}



\begin{quote}
Before the new one: say the Persian for a century, then the Persian for a term.
\end{quote}

\subsection*{You'll want to know: \fa{چهارراه}}
\textbf{\fa{چهارراه}} — \emph{chahârrâh} — \textquotedblleft{}a crossroads\textquotedblright{}.

\begin{cousinweb}
\emph{chahârrâh} is \emph{chahâr}, \textquotedblleft{}four\textquotedblright{}, with \emph{râh}, \textquotedblleft{}road\textquotedblright{}: where four roads meet. Both are inherited Persian words.
\end{cousinweb}

\begin{grammarlens}[title={What you've built}]
One more word for this chapter.
\end{grammarlens}

\subsection*{Your turn}
\begin{itemize}
\raggedright
  \item \emph{Say it:} \emph{chahârrâh}
  \item \emph{Say it:} \emph{chahârrâh}, once more
  \item \emph{Say it:} say \emph{chahârrâh}
  \item \emph{Recall:} say the Persian for a century, then the Persian for a term, then say \emph{chahârrâh} again
\end{itemize}

\begin{tcolorbox}[breakable,colback=teal!4,colframe=teal!35!black,title={Before you move on}]
What does \fa{چهارراه} mean? (\textquotedblleft{}A crossroads\textquotedblright{}.) Say it once more.
\end{tcolorbox}
\section[\texorpdfstring{\fa{محله} (mahalle)}{محله (mahalle)}]{\fa{محله} (mahalle) — a neighbourhood}

\label{lesson:FA-C266-mahalle}



\begin{quote}
Before the new one: say the Persian for a term, then the Persian for a crossroads.
\end{quote}

\subsection*{You'll want to know: \fa{محله}}
\textbf{\fa{محله}} — \emph{mahalle} — \textquotedblleft{}a neighbourhood\textquotedblright{}.

\begin{cousinweb}
\emph{mahalle} is borrowed from Arabic \emph{maḥalla}, \textquotedblleft{}quarter, place of halting\textquotedblright{}. Its Arabic root is \emph{ḥ-l-l} (Lane).
\end{cousinweb}

\begin{grammarlens}[title={What you've built}]
One more word for this chapter.
\end{grammarlens}

\subsection*{Your turn}
\begin{itemize}
\raggedright
  \item \emph{Say it:} \emph{mahalle}
  \item \emph{Say it:} \emph{mahalle}, once more
  \item \emph{Say it:} say \emph{mahalle}
  \item \emph{Recall:} say the Persian for a term, then the Persian for a crossroads, then say \emph{mahalle} again
\end{itemize}

\begin{tcolorbox}[breakable,colback=teal!4,colframe=teal!35!black,title={Before you move on}]
What does \fa{محله} mean? (\textquotedblleft{}A neighbourhood\textquotedblright{}.) Say it once more.
\end{tcolorbox}
\section[\texorpdfstring{\fa{ساختمان} (sâkhtemân)}{ساختمان (sâkhtemân)}]{\fa{ساختمان} (sâkhtemân) — a building}

\label{lesson:FA-C266-sakhteman}



\begin{quote}
Before the new one: say the Persian for a crossroads, then the Persian for a neighbourhood.
\end{quote}

\subsection*{You'll want to know: \fa{ساختمان}}
\textbf{\fa{ساختمان}} — \emph{sâkhtemân} — \textquotedblleft{}a building\textquotedblright{}.

\begin{cousinweb}
\emph{sâkhtemân} is built from \emph{sâkht}, \textquotedblleft{}made, built\textquotedblright{} (from \emph{sâkhtan}, \textquotedblleft{}to build\textquotedblright{}), and the ending \emph{-emân}: a built thing. \emph{sâkhtan} is inherited, usually traced to Middle Persian \emph{sāxtan}.
\end{cousinweb}

\begin{grammarlens}[title={What you've built}]
One more word for this chapter.
\end{grammarlens}

\subsection*{Your turn}
\begin{itemize}
\raggedright
  \item \emph{Say it:} \emph{sâkhtemân}
  \item \emph{Say it:} \emph{sâkhtemân}, once more
  \item \emph{Say it:} say \emph{sâkhtemân}
  \item \emph{Recall:} say the Persian for a crossroads, then the Persian for a neighbourhood, then say \emph{sâkhtemân} again
\end{itemize}

\begin{tcolorbox}[breakable,colback=teal!4,colframe=teal!35!black,title={Before you move on}]
What does \fa{ساختمان} mean? (\textquotedblleft{}A building\textquotedblright{}.) Say it once more.
\end{tcolorbox}
\section[\texorpdfstring{\fa{ورودی} (vorudi)}{ورودی (vorudi)}]{\fa{ورودی} (vorudi) — an entrance}

\label{lesson:FA-C266-vorudi}



\begin{quote}
Before the new one: say the Persian for a neighbourhood, then the Persian for a building.
\end{quote}

\subsection*{You'll want to know: \fa{ورودی}}
\textbf{\fa{ورودی}} — \emph{vorudi} — \textquotedblleft{}an entrance\textquotedblright{}.

\begin{cousinweb}
\emph{vorudi} is \emph{vorud}, \textquotedblleft{}entry\textquotedblright{}, with the ending \emph{-i}. \emph{vorud} is borrowed from Arabic \emph{wurūd}, \textquotedblleft{}arrival, coming to water\textquotedblright{} (root \emph{w-r-d}, Lane).
\end{cousinweb}

\begin{grammarlens}[title={What you've built}]
One more word for this chapter.
\end{grammarlens}

\subsection*{Your turn}
\begin{itemize}
\raggedright
  \item \emph{Say it:} \emph{vorudi}
  \item \emph{Say it:} \emph{vorudi}, once more
  \item \emph{Say it:} say \emph{vorudi}
  \item \emph{Recall:} say the Persian for a neighbourhood, then the Persian for a building, then say \emph{vorudi} again
\end{itemize}

\begin{tcolorbox}[breakable,colback=teal!4,colframe=teal!35!black,title={Before you move on}]
What does \fa{ورودی} mean? (\textquotedblleft{}An entrance\textquotedblright{}.) Say it once more.
\end{tcolorbox}
\section[\texorpdfstring{\fa{خروجی} (khoruji)}{خروجی (khoruji)}]{\fa{خروجی} (khoruji) — an exit}

\label{lesson:FA-C266-khoruji}



\begin{quote}
Before the new one: say the Persian for a building, then the Persian for an entrance.
\end{quote}

\subsection*{You'll want to know: \fa{خروجی}}
\textbf{\fa{خروجی}} — \emph{khoruji} — \textquotedblleft{}an exit\textquotedblright{}.

\begin{cousinweb}
\emph{khoruji} is \emph{khoruj}, \textquotedblleft{}going out\textquotedblright{}, with the ending \emph{-i}. \emph{khoruj} is borrowed from Arabic \emph{ḫurūj}, \textquotedblleft{}going out\textquotedblright{} (root \emph{ḫ-r-j}, Lane).
\end{cousinweb}

\begin{grammarlens}[title={What you've built}]
One more word for this chapter.
\end{grammarlens}

\subsection*{Your turn}
\begin{itemize}
\raggedright
  \item \emph{Say it:} \emph{khoruji}
  \item \emph{Say it:} \emph{khoruji}, once more
  \item \emph{Say it:} say \emph{khoruji}
  \item \emph{Recall:} say the Persian for a building, then the Persian for an entrance, then say \emph{khoruji} again
\end{itemize}

\begin{tcolorbox}[breakable,colback=teal!4,colframe=teal!35!black,title={Before you move on}]
What does \fa{خروجی} mean? (\textquotedblleft{}An exit\textquotedblright{}.) Say it once more.
\end{tcolorbox}
